% !TeX root = ../main.tex
% The keywords are taken from \Keywords in config.tex
\begin{abstractpage}

\begin{guidance}
The abstract is the most widely read part of your thesis, and often the only one that
appears in your university's repository. Write it last, as a self-contained text of
about 200--300 words (check the limit of your university) that covers, in this order:
the \emph{context} and the problem, the \emph{objective}, the \emph{methods}, the main
\emph{results} (with numbers when possible) and the \emph{conclusions}. Avoid
citations, undefined acronyms, equations and references to figures or tables.
\end{guidance}

This document is a \LaTeX{} template for Master's theses written in English at Spanish
public universities, and at the same time it is its own user guide. It provides a
complete and configurable structure, from the cover page and the declarations commonly
required by Spanish regulations to the chapters of a research-oriented thesis, the
references and the appendices. All the personal and institutional data are defined in a
single configuration file, so the same template can be used at any university by changing
the university name, logos, colours and a few options. The template supports digital and
printed versions, two typefaces, a second abstract in Spanish or in any co-official
language, and guidance boxes that disappear automatically in the final version. It uses
only the packages of the original template, so it is light and fast to compile. Replace this text with your own abstract.

\end{abstractpage}
